% Part: counterfactuals
% Chapter: introduction
% Section: paradoxes-material

\documentclass[../../../include/open-logic-section]{subfiles}

\begin{document}

\olfileid{cnt}{int}{par}

\olsection{Paradoks Kondisional Material}

Salah siji wong kapisan kang ngritik panggunaan $!A \lif !B$ kanggo
nyimbolake pratelan basa Inggris ``if \dots then \dots'' yaiku
C.~I. Lewis. Lewis ngritik panggunaan kondisional material ing
\emph{Principia Mathematica} anggitane Whitehead lan Russell, kang
maca $\lif$ minangka ``implies'' (``ngimplikasikake''). Lewis kanthi
bener ngandharake yen $\lif$ tegese ``implies'', mula saben
proposisi~$p$ kang luput ngimplikasikake yen $p$ ngimplikasikake~$q$,
amarga $p \lif (p \lif q)$ bener yen~$p$ luput; lan saben proposisi~$q$
kang bener ngimplikasikake yen $p$ ngimplikasikake~$q$, amarga
$q \lif (p \lif q)$ bener yen $q$~bener.

Mesthi wae para ahli logika mangerteni yen \emph{implikasi}, yaiku
konsekuensi logis, iku relasi antarane !!{formula}s utawa pratelan,
dudu konektif. Mula kanggo nyingkiri kebingungan, kita aja maca
$\lif$ minangka ``implies''.\footnote{Maca ``$\lif$'' minangka
  ``implies'' isih lumrah ing antarane matematikawan lan ilmuwan
  komputer, sanadyan para filsuf ngupaya nyingkiri kebingungan kang
  dituduhake Lewis kanthi maca minangka ``only if'' (``mung yen'').}
Anggere kita ora maca mangkono, prakara kang dikuwatirake Lewis
ora tuwuh: $p$ ora ``ngimplikasikake'' $q$, sanajan kita nganggep
$p$ minangka wakil ukara basa Inggris kang luput. Kanggo nemtokake
apa $p \Entails q$, kita kudu nimbang \emph{kabeh} !!{valuation}s,
lan $p \Entails/ q$ sanajan kita nggunakake $p$ kanggo nyimbolake
ukara kang kebeneran luput.

Nanging isih ana kang aneh ing pratelan ``if \dots then\dots''
kaya tuladhane Lewis iki:
\begin{quote}
Yen rembulan digawe saka keju ijo, mula $2+2=4$.
\end{quote}
lan ing inferensi iki:
\begin{quote}
  Rembulan ora digawe saka keju ijo. Mula, yen rembulan digawe
  saka keju ijo, mula $2+2=4$.

  $2+2 = 4$. Mula, yen rembulan digawe
  saka keju ijo, mula $2+2=4$.
\end{quote}
Nanging yen ``if \dots then \dots'' mung ateges $\lif$, ukara
kasebut bener tanpa masalah, lan inferensi-inferensine sah tanpa
masalah.

Tuladha liyane gegayutan karo tautologi $(!A \lif !B) \lor (!B \lif
!A)$. Iki nuduhake yen kita njupuk rong ukara indikatif~$S$ lan $T$
saka koran kanthi acak, ukara ``Yen $S$ mula $T$, utawa yen $T$
mula~$S$'' mesthine bener.

\end{document}
